% ======================================================================================
% Nom     : reports/latex/sections/00_glossaire.tex
% Rôle    : Liste des sigles et glossaire des termes employés dans le rapport.
% Auteur  : Maxime BRONNY
% Version : V1
% Cadre   : Réalisé dans le cadre du cours Techniques d'apprentissage artificiel
%           Master 1 Informatique Big Data
% Usage   : Fichier inclus dans main.tex lors de la compilation du rapport.
% ======================================================================================

% ============================================================================
% Liste des sigles et glossaire - termes et acronymes réellement employés dans
% le rapport, classés par ordre alphabétique (terme à gauche, définition à
% droite). Page liminaire référencée dans le sommaire.
%
% Mise en page : longtable à deux colonnes sans filets (terme en gras à gauche,
% définition alignée à droite), qui se répartit proprement sur plusieurs pages
% et réserve toujours la hauteur de la ligne (pas de chevauchement).
% ============================================================================
\clearpage
\section*{Liste des sigles et glossaire}
\addcontentsline{toc}{section}{Liste des sigles et glossaire}

\begingroup
\setlength{\extrarowheight}{0.35em}   % respiration verticale entre les entrées
\begin{longtable}{@{}>{\raggedright\bfseries}p{3.2cm}@{\hspace{0.4cm}}>{\raggedright\arraybackslash}p{\dimexpr\textwidth-3.7cm\relax}@{}}

Accuracy & Exactitude : proportion de prédictions correctes, toutes classes
  confondues, $(TP+TN)/N$. Peu informative seule sur des données déséquilibrées. \\

Apprentissage supervisé & Famille de méthodes qui apprennent à prédire une
  variable cible à partir d'exemples étiquetés. \\

Arbre de décision & Modèle qui partitionne récursivement les données par des
  tests binaires « variable $\leq$ seuil » et prédit, dans chaque feuille, la
  proportion de défauts observée (voir CART). \\

AUC & \emph{Area Under the Curve} : aire sous la courbe ROC ; probabilité qu'un
  défaut tiré au hasard reçoive un score plus élevé qu'un non-défaut
  ($0{,}5$~=~hasard, $1$~=~classement parfait). \\

CART & \emph{Classification And Regression Tree} : algorithme d'arbre de décision
  fondé sur des séparations binaires successives des données. \\

Classification binaire & Tâche supervisée où la cible ne prend que deux
  valeurs ; ici, $0$~=~remboursement, $1$~=~défaut de paiement. \\

Courbe ROC & Courbe du taux de vrais positifs en fonction du taux de faux
  positifs lorsque le seuil de décision varie (voir ROC, AUC). \\

Dataset & Jeu de données : ensemble des observations disponibles ; ici,
  32\,581 dossiers de crédit décrits par 11 variables et une cible. \\

Descente de gradient & Algorithme d'optimisation itératif qui ajuste les
  paramètres dans la direction qui fait le plus baisser la fonction de coût. \\

Encodage & Conversion des variables catégorielles en nombres ; nous utilisons un
  \emph{label encoding} à correspondance fixe. \\

F1-score & Moyenne harmonique de la précision et du rappel ; critère de
  sélection des hyperparamètres dans ce projet. \\

Faux négatif (FN) & Défaut réel non détecté (crédit accordé à un client qui ne
  remboursera pas) : l'erreur la plus coûteuse pour la banque. \\

Faux positif (FP) & Bon client classé à tort comme défaut (crédit refusé sans
  raison réelle). \\

FN & \emph{Faux négatif} : voir Faux négatif. \\

FP & \emph{Faux positif} : voir Faux positif. \\

From scratch & Implémenté à partir des opérations de base (NumPy), sans
  bibliothèque d'apprentissage automatique. \\

Fuite de données & \emph{Data leakage} : utilisation, même indirecte,
  d'informations de validation ou de test pendant la conception du modèle ;
  elle conduit à surestimer les performances. \\

Généralisation & Capacité d'un modèle à rester performant sur des données jamais
  vues ; c'est ce que mesure l'ensemble de test. \\

Hyperparamètre & Réglage fixé avant l'apprentissage (profondeur de l'arbre,
  nombre de voisins $k$, seuil de décision), choisi ici sur l'ensemble de
  validation. \\

Imputation & Remplacement des valeurs manquantes par une valeur calculée ; ici,
  la médiane de l'ensemble d'entraînement. \\

Impureté de Gini & Critère de qualité d'une division de l'arbre : $G = 2p(1-p)$
  pour un groupe contenant une proportion $p$ de défauts. \\

Jeu d'entraînement & Sous-ensemble (70~\%) servant à apprendre les paramètres
  des modèles. \\

Jeu de test & Sous-ensemble (15~\%) utilisé une seule fois, à la fin, pour
  mesurer la performance de généralisation. \\

Jeu de validation & Sous-ensemble (15~\%) servant à choisir les
  hyperparamètres. \\

k-NN & \emph{k plus proches voisins} (\emph{k-Nearest Neighbors}) : méthode qui
  classe un exemple par vote des $k$ exemples d'entraînement les plus proches au
  sens de la distance euclidienne. \\

Log-loss & Entropie croisée : fonction de coût de la régression logistique,
  mesurant l'écart entre probabilités prédites et étiquettes réelles. \\

Matrice de confusion & Tableau croisant classes réelles et classes prédites
  (TP, TN, FP, FN). \\

Modèle & Fonction apprise des données qui associe une prédiction à un vecteur de
  variables explicatives. \\

Normalisation & Mise à l'échelle de chaque variable (score $z$) : soustraction
  de la moyenne puis division par l'écart-type, calculés sur l'ensemble
  d'entraînement. \\

NumPy & Bibliothèque Python de calcul numérique (tableaux, algèbre linéaire) ;
  base de toutes les implémentations \emph{from scratch} du projet. \\

Précision & Parmi les défauts prédits, proportion de défauts réels :
  $TP/(TP+FP)$. \\

Prétraitement & Transformations appliquées aux données brutes avant
  l'apprentissage : nettoyage, encodage, imputation, normalisation. \\

Python & Langage de programmation utilisé pour l'ensemble du projet. \\

Rappel & Parmi les défauts réels, proportion détectée : $TP/(TP+FN)$. \\

Régression logistique & Modèle linéaire qui estime la probabilité de défaut en
  passant le score $\mathbf{w}\cdot\mathbf{x}+b$ dans la sigmoïde. \\

ROC & \emph{Receiver Operating Characteristic} : voir Courbe ROC. \\

Scikit-learn & Bibliothèque Python d'apprentissage automatique, utilisée ici
  uniquement pour vérifier la justesse des implémentations \emph{from scratch}. \\

Seuil de décision & Valeur au-delà de laquelle la probabilité prédite est
  convertie en classe « défaut ». \\

Sigmoïde & Fonction $\sigma(z) = 1/(1+e^{-z})$, qui transforme un score réel en
  probabilité comprise entre $0$ et $1$. \\

Stratification & Découpage effectué classe par classe afin que chaque ensemble
  conserve la même proportion de défauts ($21{,}8\,\%$). \\

Surapprentissage & \emph{Overfitting} : situation où le modèle colle au bruit de
  l'entraînement et généralise mal à de nouvelles données. \\

TN & \emph{True Negative} (vrai négatif) : remboursement correctement
  identifié. \\

TP & \emph{True Positive} (vrai positif) : défaut correctement détecté. \\

Variable cible & Valeur à prédire (\texttt{loan\_status}) : $1$~=~défaut,
  $0$~=~remboursement. \\

Variable explicative & Entrée du modèle (âge, revenu, montant, taux, etc.) à
  partir de laquelle il prédit la cible. \\

\end{longtable}
\endgroup
